% Paper 5, §5 — Reasoning benchmarks: where saturation does and doesn't happen
% Anchors: kb/notes/evaluation/reasoning-benchmarks.md
%          kb/excerpts/{glazer2024-frontiermath,rein2023-gpqa,phan2025-hle}.md
%          kb/excerpts/hendrycks2021-mmlu.md (saturation backstory)

\section{Reasoning benchmarks: where saturation does and doesn't happen}
\label{sec:reasoning-benchmarks}

The lineage in \cref{sec:knowledge-benchmarks} measured what a model
\emph{retrieves}; the lineage in \cref{sec:contamination} measured how
much the test set has \emph{leaked} into the training distribution.
This section measures something distinct: \emph{multi-step reasoning}
on items where the answer is not directly recoverable from a single
fact or a single web hit. The 2024--2026 reasoning-benchmark surface
splits cleanly into two regimes. \emph{Contest-style} math
(AIME~2024, MATH-500) and \emph{graduate-question Q\&A} (\glsterm{gpqa-diamond-2}{GPQA Diamond})
have saturated at the frontier --- the very regime that Paper~3's
\glsterm{rlvr-3}{RLVR}/\glsterm{grpo-2}{GRPO} reasoning-RL lineage was trained against, where the
verifiable-reward signal is dense and the problem distribution has
been seen by every lab's training pipeline. \emph{Research-style}
math (\glsterm{frontiermath}{FrontierMath} \citep{glazer2024-frontiermath}) and
\emph{broad-coverage frontier-difficulty exams} (Humanity's Last
Exam, \glsterm{hle}{HLE} \citep{phan2025-hle}) have not. The thesis of this section
is that the gap between these two regimes is the gap between
in-distribution gains from reasoning-RL and out-of-distribution
generalization, and the field's 2026 picture is that the latter is
much weaker than the former.

The shared formal frame is the reasoning-benchmark instance
$\mathcal{B} = (\rho, A, V)$%
\footnote{KB note: \texttt{kb/notes/evaluation/reasoning-benchmarks.md\#1-the-reasoning-benchmark-protocol}.}
--- a problem statement $\rho$, an answer-format constraint $A$
(numeric integer, single letter, parsed symbolic expression), and
a verifier $V(\hat{a}, a^\star) \in \{0, 1\}$ on the model's chosen
answer. There is no environment, no trajectory, no end-state ---
$V$ is exact-match or a domain-specific parser, by construction tight
enough to resist lucky guesses and automated enough to run at scale
without \glsterm{llm-as-judge}{LLM-as-judge}. What separates the four reasoning benchmarks
below is the difficulty distribution of $\rho$ and the human-effort
calibration that fixes the \glsterm{saturation}{saturation} timeline.

\subsection{The contest-math regime: AIME and MATH-500}
\label{sec:reasoning-aime}

Contest-style math is the regime where reasoning-RL has worked best.
The two canonical benchmarks both predate the post-R1 reasoning
generation by several years. MATH \deepencite{hendrycks2021-math §1}
ships $12{,}500$ competition-style problems, most middle- and
high-school level, with numeric or LaTeX-rendered free-form answers
checked by symbolic match%
\footnote{KB note: \texttt{kb/notes/evaluation/reasoning-benchmarks.md\#3-1-the-post-2024-reasoning-benchmark-map}.};
\textbf{MATH-500} is the $500$-item subset that the reasoning literature
adopted as a fast-iterating slice. AIME~2024 is the $30$-problem
American Invitational Mathematics Exam set with integer answers in
$[0, 999]$, three hours of human work for a strong competitor, and
pre-2025 launch numbers in the $\sim 50\%$ range for early
reasoning-tuned models \deepencite{aime-2024 §1}.

The 2025 reasoning-RL generation moved both numbers decisively. As
documented in Paper~3's reasoning-RL section,
DeepSeek-R1's RL-only training
pipeline reports an AIME~2024 jump from a base-model $\sim 16\%$ to
a final $\sim 80\%$ pass$@1$
\deepencite{deepseek-r1 §3}, with \glsterm{self-consistency-2}{self-consistency} over $64$ samples
pushing the number above $90\%$. By 2026-Q1 frontier reasoning models
report AIME~2024 in the $83$--$100\%$ band%
\footnote{KB note: \texttt{kb/notes/evaluation/reasoning-benchmarks.md\#3-1-the-post-2024-reasoning-benchmark-map}, \emph{\glsterm{saturation}{Saturation} status} column.},
with the upper end pinned to perfect on the $30$-problem set --- the
benchmark has effectively been saturated. MATH-500 follows the same
arc; published reasoning-model numbers above $95\%$ are common.

Two factors compose the AIME-2024 \glsterm{saturation}{saturation} signal.
\textbf{Reasoning-RL on verifiable rewards.} AIME's integer-answer
verifier is the canonical \glsterm{rlvr-3}{RLVR} target: the reward is $\{0, 1\}$ on
exact-match against the gold integer, no judge model, no
preference labels, no \glsterm{value}{value} head. Paper~3's \glsterm{grpo-2}{GRPO}/\glsterm{rlvr-3}{RLVR} lineage
(\deepencite{shao2024 §3}, \deepencite{dapo2025 §2}) is built around
exactly this signal. In-distribution gains on the trained problem
distribution are large.
\textbf{Plausible contamination.} AIME~2024 was published in
February~2024; the problems and worked solutions appeared in public
math forums within days. Pre-training corpora and reasoning-RL
trajectory data assembled in 2024-Q3 onward almost certainly contain
the items, often with worked solutions inline. \cref{sec:contamination}
sets up the AIME~2024 timeline as the case study; here the relevant
implication is that the AIME jump cannot be cleanly attributed to
RL-induced reasoning generalization versus exposure during training.

\subsection{The graduate-Q\&A regime: GPQA Diamond saturates}
\label{sec:reasoning-gpqa}

\glsterm{gpqa-diamond-2}{GPQA Diamond} \citep{rein2023-gpqa} sits in a different point of the
reasoning surface than AIME --- $198$ four-way multiple-choice
questions in graduate-level biology, physics, and chemistry, where
the difficulty comes from domain depth rather than from a tower of
algebraic manipulation. \cref{sec:knowledge-gpqa} treats the
benchmark's launch construction; the reasoning-relevant fact is
the \glsterm{saturation}{saturation} rate%
\footnote{KB excerpt: \texttt{kb/excerpts/rein2023-gpqa.md\#sec-ai-baseline}.}.
At launch (November~2023) the strongest GPT-4-based baseline
reached $39\%$ \citep[\S{}sec-ai-baseline]{rein2023-gpqa}, just above
the $34\%$ non-expert-with-web ceiling and well below the
$\sim 65\%$ PhD-expert baseline%
\footnote{KB excerpt: \texttt{kb/excerpts/rein2023-gpqa.md\#sec-baseline}.}.
By 2026-Q1, frontier reasoning models score above $90\%$ on Diamond%
\footnote{KB note: \texttt{kb/notes/evaluation/reasoning-benchmarks.md\#3-1-the-post-2024-reasoning-benchmark-map}; cross-referenced from public leaderboards.}
--- $\sim 25$ percentage points above the human PhD baseline, on a
benchmark constructed to resist web search at the $34\%$ ceiling. The
$39\% \to 90\%+$ jump in $\sim 2.5$ years exceeds the AIME
trajectory in absolute capability gain.

GPQA's own framing places this \glsterm{saturation}{saturation} in tension with its
intended use. The abstract names the design target as scalable
oversight%
\footnote{KB excerpt: \texttt{kb/excerpts/rein2023-gpqa.md\#sec-why}.}
\citep[\S{}sec-why]{rein2023-gpqa}:

\begin{quote}\itshape
we need to develop \glsterm{scalable-oversight}{scalable oversight} methods that enable humans to
supervise their outputs, which may be difficult even if the
supervisors are themselves skilled and knowledgeable.
\end{quote}

The reasoning-relevant reading is that GPQA was \emph{built to be
hard} for skilled humans even with web access, on the assumption
that the items would remain hard for AI for some years. That
assumption did not survive 2024--2025: the same RL-on-verifiable-
rewards lineage that closed AIME also closed \glsterm{gpqa-diamond-2}{GPQA Diamond}, despite
the \glsterm{google-proof-construction}{Google-proof construction}.

\begin{intuition}
Contest-style math (AIME, MATH-500) and graduate-Q\&A
(\glsterm{gpqa-diamond-2}{GPQA Diamond}) are different problem distributions but the same
\emph{distributional kind} from the reasoning-RL training pipeline's
point of view: closed-form items with a single canonical answer, a
finite gold set that is in-or-near-distribution to the RL training
mixture, and a verifier $V(\hat{a}, a^\star)$ tight enough to feed a
binary reward. \glsterm{rlvr-3}{RLVR} on this kind of distribution generalizes
\emph{within} the distribution in the way Paper~3's reasoning-RL
results document. What the \glsterm{frontiermath}{FrontierMath} gap below shows is that this
generalization is much weaker out of distribution: the same RL
pipeline does not transfer to research-mathematician problems that
demand multi-day proof construction, even though the verifier
discipline ($V$ as integer-or-symbolic match) is identical.
(Returning to canonical form: $\mathcal{B} = (\rho, A, V)$ has the
same shape across all four; what changes is only the difficulty
distribution of $\rho$ and how much of that distribution the RL
training mixture has covered.)
\end{intuition}

\subsection{FrontierMath: research-style math is the OOD bar}
\label{sec:reasoning-frontiermath}

\citet{glazer2024-frontiermath} construct \glsterm{frontiermath}{FrontierMath} as the
deliberate out-of-distribution counterweight%
\footnote{KB excerpt: \texttt{kb/excerpts/glazer2024-frontiermath.md\#abstract}.}
\citep[\S{}abstract]{glazer2024-frontiermath}:

\begin{quote}\itshape
We introduce \glsterm{frontiermath}{FrontierMath}, a benchmark of hundreds of original,
exceptionally challenging mathematics problems crafted and vetted by
expert mathematicians. The questions cover most major branches of
modern mathematics --- from computationally intensive problems in
number theory and real analysis to abstract questions in algebraic
geometry and category theory.
\end{quote}

Three construction choices distinguish \glsterm{frontiermath}{FrontierMath} from the
contest-math line and bear directly on the \glsterm{saturation}{saturation} question.
\textbf{Unpublished problems.} The problem pool is constructed by
$60+$ working mathematicians (verified PhD or equivalent), each
authoring items in their domain; problems are cross-vetted by a
second mathematician for solvability and uniqueness of answer
\citep[\S{}sec-verification]{glazer2024-frontiermath}. The set is
private; items have not appeared in public math forums prior to
benchmark release, which closes the dominant contamination channel
that AIME~2024 left open.
\textbf{Single-canonical-answer parsing.} Each problem has a
single integer, polynomial, or closed-form expression as its gold
answer, parseable by an automated checker without \glsterm{llm-as-judge}{LLM-as-judge}%
\footnote{KB excerpt: \texttt{kb/excerpts/glazer2024-frontiermath.md\#sec-verification}.}.
Problems requiring proof or open-ended argument are excluded, which
is what enables \glsterm{rlvr-3}{RLVR}-style reward computation on a benchmark that
otherwise looks like research math.
\textbf{Multi-day-effort calibration.} The difficulty discipline is
explicit \citep[\S{}sec-difficulty]{glazer2024-frontiermath}:
``Solving a typical problem requires multiple hours of effort from a
researcher in the relevant branch of mathematics, and for the upper
end questions, multiple days.'' Where MATH problems take a competent
undergraduate $\sim 30$~min and AIME problems take a strong
high-schooler $\sim 10$~min, \glsterm{frontiermath}{FrontierMath} items are calibrated to
specialist research effort%
\footnote{KB excerpt: \texttt{kb/excerpts/glazer2024-frontiermath.md\#sec-difficulty}.}.

The headline launch number is the load-bearing one
\citep[\S{}abstract,\S{}sec-launch-result]{glazer2024-frontiermath}:

\begin{quote}\itshape
Current state-of-the-art AI models solve under 2\% of problems,
revealing a vast gap between AI capabilities and the prowess of the
mathematical community.
\end{quote}

The under-$2\%$ floor was reported on models that were already at
$80$--$95\%$ on \glsterm{mmlu}{MMLU} and post-\glsterm{saturation}{saturation} on AIME~2024. The gap is
the OOD bar: the same model class that has saturated the
contest-style distribution scores in the single digits on the
research-style one. By the close of 2025, frontier \glsterm{reasoning-model}{reasoning model}
reports on \glsterm{frontiermath}{FrontierMath} had moved into the low double digits ---
single-to-low-double-digit accuracy is the 2026-frontier band on the
benchmark%
\footnote{KB note: \texttt{kb/notes/evaluation/reasoning-benchmarks.md\#3-1-the-post-2024-reasoning-benchmark-map}, \emph{actively discriminative} status as of 2026-05.}
--- two orders of magnitude below contest-math \glsterm{saturation}{saturation}. The gap
is the empirical content of the in-distribution-vs-OOD claim.

\begin{contradiction}
\textbf{Is \glsterm{frontiermath}{FrontierMath} saturating too?} The benchmark is held by a
single group (Epoch AI) and graded against a private answer \glsterm{key}{key} with
a domain-specific symbolic checker. The construction is
contamination-resistant by design, but it is also single-source: the
field cannot independently replicate either the problem set or the
grading infrastructure, so reported \glsterm{frontiermath}{FrontierMath} scores cannot be
audited by a third party in the way \glsterm{gpqa-diamond-2}{GPQA Diamond} and AIME~2024 can.
Reasoning-model \glsterm{frontiermath}{FrontierMath} numbers reported in late 2025 and early
2026 vary substantially across vendor announcements, partly because
the evaluation protocol --- which problem subset, which compute
budget, which sampling strategy --- is not fixed by a public spec.
The contradiction is between (i) the original $<\!2\%$ launch
number, which is well-bounded and externally verifiable in the loose
sense that ``almost everyone scored close to zero'' is hard to fake,
and (ii) the post-2025 numbers in the low double digits, which are
single-source claims on a single-source benchmark. The field has not
shipped a replicate \glsterm{frontiermath}{FrontierMath}-class benchmark with an open
problem-construction protocol; until it does, \glsterm{frontiermath}{FrontierMath}
\glsterm{saturation}{saturation} rate is what its custodians say it is.
\end{contradiction}

\subsection{HLE as the broad-coverage OOD frontier}
\label{sec:reasoning-hle}

\citet{phan2025-hle} extend the \glsterm{frontiermath}{FrontierMath} move from one
discipline to many. \glsterm{hle}{Humanity's Last Exam} ships $2{,}500$ questions
across mathematics, the humanities, and the natural sciences,
mixing multiple-choice and short-answer items, with a verifier that
is choice-equality or string-equality-after-normalization --- not
\glsterm{llm-as-judge}{LLM-as-judge}%
\footnote{KB excerpt: \texttt{kb/excerpts/phan2025-hle.md\#sec-format}.}.
The motivation paragraph names the gap that \glsterm{hle}{HLE} fills
\citep[\S{}abstract]{phan2025-hle}:

\begin{quote}\itshape
LLMs now achieve over 90\% accuracy on popular benchmarks like \glsterm{mmlu}{MMLU},
limiting informed measurement of state-of-the-art \glsterm{llm}{LLM} capabilities.
\end{quote}

The construction inherits the GPQA / \glsterm{frontiermath}{FrontierMath} verifier discipline
\citep[\S{}sec-format]{phan2025-hle}:

\begin{quote}\itshape
a known solution that is unambiguous and easily verifiable, but
cannot be quickly answered via internet retrieval.
\end{quote}

The capability gap at publication is the property the design was
reaching for \citep[\S{}sec-capability-gap]{phan2025-hle}: ``State-of-
the-art LLMs demonstrate low accuracy and calibration on \glsterm{hle}{HLE},
highlighting a significant gap between current \glsterm{llm}{LLM} capabilities and
the expert human frontier on closed-ended academic questions.'' By
the close of the Phase~1 sweep window in early 2026, top model
scores on \glsterm{hle}{HLE} were in the $\sim 65\%$ band on a benchmark designed
to resist web retrieval and frontier reasoning at construction time%
\footnote{KB note: \texttt{kb/notes/evaluation/reasoning-benchmarks.md\#5-frontier-and-open-questions-as-of-2026-05}.}
--- substantially below \glsterm{mmlu}{MMLU} \glsterm{saturation}{saturation} but well above \glsterm{frontiermath}{FrontierMath}.
The position of \glsterm{hle}{HLE} on the \glsterm{saturation}{saturation} curve is what makes it the
useful broad-coverage OOD instrument in 2026: it is hard enough that
contest-math RL transfer does not close it, and broad enough that
single-domain RL transfer would not close it either, but not so far
out of distribution as to floor at single digits.

\subsection{What the split tells us about reasoning-RL}
\label{sec:reasoning-rl-implications}

The four-benchmark surface produces a single load-bearing claim
about Paper~3's reasoning-RL lineage. Reasoning-RL on verifiable
rewards (Paper~3 §\glsterm{rlvr-3}{RLVR}-and-\glsterm{grpo-2}{GRPO}) delivers large gains on
problem distributions covered by the RL training mixture --- the
$16\% \to 80\%+$ AIME~2024 jump documented in
\citet[\S{}3]{deepseek-r1}, and the parallel \glsterm{saturation}{saturation} of MATH-500
and \glsterm{gpqa-diamond-2}{GPQA Diamond}, are real and they generalize within the contest /
graduate-Q\&A distributional kind. The same lineage delivers
single-digit-to-low-double-digit accuracy on \glsterm{frontiermath}{FrontierMath}, where
the problem distribution sits two human-effort orders of magnitude
above the RL training distribution and outside the public
contest-and-PhD-exam corpus that the RL trajectory data is built
from. The RL-induced gains do not transfer cleanly into this regime.

This is consistent with a moderate version of the
in-distribution-RL claim: \glsterm{rlvr-3}{RLVR}-on-verifiable-rewards is doing
real reasoning-\glsterm{circuit}{circuit} work within its training distribution
(Paper~3 §\glsterm{rlvr-3}{RLVR}-generalization-vs-memorization), but
its OOD generalization is much weaker than the AIME-and-GPQA
trajectories alone would suggest. The $80\%+$ AIME \glsterm{saturation}{saturation} does
not predict $80\%+$ on \glsterm{frontiermath}{FrontierMath}; it predicts $1$--$2\%$ on
\glsterm{frontiermath}{FrontierMath}, which is what the benchmarks observe. The frontier
field's read of the gap by mid-2026 is that further AIME-style
\glsterm{saturation}{saturation} is no longer informative for capability measurement,
and that the next round of measurement infrastructure has to
target \glsterm{frontiermath}{FrontierMath}-and-\glsterm{hle}{HLE}-class items where the OOD gap is the
construct \deepencite{glazer2024-frontiermath §sec-saturation}.

The cross-cite to \cref{sec:contamination} is sharper here than
elsewhere in the paper. \cref{sec:contamination} establishes
AIME~2024 as the canonical contamination case study: the items and
worked solutions are public, were almost certainly absorbed into
2024-Q3+ pre-training corpora and reasoning-RL trajectory data, and
the line between ``the model knows the problem distribution'' and
``the model has seen the items'' is fuzzy and field-disputed. The
implication for AIME-jump interpretation is that the reasoning
community cannot cleanly attribute the $16\% \to 80\%+$ gain to
RL-induced reasoning generalization versus contamination-aided
recall, because both are present and the experimental design does
not separate them. \glsterm{frontiermath}{FrontierMath} is the contamination defense the
contest line lacks; the empirical OOD gap is the part of the
reasoning-RL story that the contamination defense leaves open.

\subsection{Forward-link: from reasoning to agentic}
\label{sec:reasoning-forward}

\cref{sec:agentic-benchmarks} shifts the eval surface from single-
response reasoning to multi-step trajectories in an environment ---
SWE-Bench, OSWorld, GAIA, $\tau$-Bench. The realism-vs-measurability
tradeoff that organizes that section is the agentic counterpart of
the in-distribution-vs-OOD tradeoff that organizes this one: each
\glsterm{agentic-benchmark}{agentic benchmark} picks a corner of how much of real deployment
the eval reproduces, and pays the corresponding price in grading
reliability. \cref{sec:contradictions-live} catalogues the
\emph{is \glsterm{frontiermath}{FrontierMath} saturating too?} contradiction surfaced above
alongside the field's other live disputes; the closing
question --- \emph{what would close it?} --- is a public,
replicable research-style math benchmark whose problem-construction
and grading infrastructure are open enough to audit, and which the
field does not yet have.
